% element: 10-s052:e04
\BeleskeID{10-s052}
U inverznom aktivnom režimu fizički kolektor preuzima ulogu emitera, a fizički emitor ulogu kolektora. Zato najpre razlikujemo indekse sa oznakom „novi“ od fizičkih oznaka na simbolu.

$\beta_R$ je strujno pojačanje u inverznom aktivnom režimu i znatno je manje od uobičajenog direktnog pojačanja; u prikazanoj analizi opseg je $0.1<\beta_R<5$. Kada se vratimo fizičkim referentnim smerovima, struje dobijaju negativne znakove prikazane u donjem paru izraza.

Isto razlikovanje indeksa važno je za znak napona inverznog zasićenja. $V_{CESI}$ je napon zasićenja u tom režimu. Uvećani deo karakteristike pomaže da se vidi mali napon prelaza u području suprotnog polariteta.
